\section{Introduction}
\label{sec:intro}
The investment value of a return forecast depends on the positions it produces and the cost of changing those positions. A model with a small statistical signal can generate frequent trades, while a strategy with lower volatility can simply spend less time invested. We study these issues through a daily SPY timing strategy whose inputs, estimation, executions, and accounting are specified together.

The research question is whether recent price trends and volatility, combined through a regularized linear model, improve the allocation between SPY and cash. The feature set consists of returns over one, five, and twenty sessions, the distance from a twenty-session moving average, and recent return volatility. This specification allows short-term continuation and reversal to compete in estimation. Trends can persist when information and portfolio adjustments are gradual; temporary price pressure can instead reverse. Volatility supplies an observable measure of the environment in which those effects occur.

Ridge regression is a suitable estimator for this design because several inputs summarize overlapping price history. Its $\ell_2$ penalty stabilizes their joint coefficients while retaining all five predictors \citep{hoerl1970ridge}. The strategy is updated monthly using a fixed-length rolling window. A small validation search selects the penalty, and a separate historical period evaluates the selected trading rule. This connects model fitting to an investor's decision rather than treating forecast accuracy as the sole objective.

The distinction between statistical and economic performance is central in financial prediction. \citet{chinco2019sparse} combine regularized forecasts with an explicit trading rule at an intraday horizon; \citet{kelly2024virtue} examine how shrinkage and model complexity affect return forecasts and timing portfolios. Our study applies these design principles to a small daily model. A rolling historical mean provides a forecast reference, following the benchmark discipline emphasized by \citet{campbell2008predicting}. SPY buy-and-hold and a fixed moving-average strategy provide the investment comparisons.

Three questions organize the analysis. Does the model improve prediction relative to the historical mean? Does the selected trading rule improve performance after costs? How much of its lower risk reflects reduced equity exposure? Answering them requires a consistent chain from the information available at the decision close to the return earned after the next opening execution.

The final evaluation covers 1,002 holding intervals in 2022--2025. At 5 basis points per one-way trade, Ridge timing earns \RidgeCAGR{} CAGR and a \RidgeSharpe{} Sharpe ratio, below buy-and-hold's \BHCAGR{} and \BHSharpe{}. Its maximum drawdown is smaller, but it is invested in SPY during only \RidgeExposure{} of intervals. Prediction error is slightly higher than the historical-mean benchmark. The economic explanation is therefore not simply that costs destroy a strong forecast: gross risk-adjusted performance is already close to the passive alternative, and frequent full-position switches then reduce net returns.

The report follows that explanation. Sections~\ref{sec:data}--\ref{sec:accounting} establish the information set, estimation procedure, and portfolio accounting. Section~\ref{sec:results} links prediction diagnostics to wealth, costs, exposure, and annual outcomes. The final sections discuss the investment implications, limitations, and reproducibility of the study.
